\documentclass[11pt,a4paper]{article}
\usepackage[T1]{fontenc}
\usepackage[utf8]{inputenc}
\usepackage{geometry}
\usepackage{parskip}
\usepackage{hyperref}
\geometry{top=25mm, bottom=25mm, left=25mm, right=25mm}
\hypersetup{colorlinks=true, urlcolor=blue, linkcolor=black}
\begin{document}

Maryam Jalili \\
Department of Mathematics \\
Ne. C., Islamic Azad University \\
Neyshabur, Iran \\
\href{mailto:mrmjlili@iau.ac.ir}{mrmjlili@iau.ac.ir}

15 October 2026

Editor-in-Chief \\
\emph{Journal of Applied Mathematics and Computing} \\
Springer

Dear Editor,

I am pleased to submit my manuscript, \textbf{``A Diagnostic Framework for Detecting False Fractionalization in Physics-Informed Neural Networks: A Three-Way Comparison of Fractional Advection--Dispersion Residuals,'} for consideration for publication in the \emph{Journal of Applied Mathematics and Computing}.

\textbf{Motivation and contribution.} Physics-informed neural networks (PINNs) are increasingly deployed for fractional partial differential equations, yet the field lacks a systematic method to verify that a chosen residual discretization genuinely preserves the non-local memory of the underlying fractional operator. This paper addresses that gap through three contributions:

\begin{enumerate}
\item A \emph{three-test diagnostic protocol}---kernel-presence, $\Delta t$-refinement, and effective-time---that classifies a PINN residual as either a genuine fractional discretization or a local surrogate.
\item A \emph{local-limit lemma} for the naive two-point finite-difference approximation of the Atangana--Baleanu--Caputo derivative, with an explicit first-order error bound and a discussion of its scope relative to the singular step-inlet data used in the experiments.
\item The first \emph{three-way comparison} of PINN residuals guided by an explicit diagnostic protocol, at two fractional orders ($\alpha = 0.8$ and $\alpha = 0.5$).
\end{enumerate}

\textbf{Reproducibility.} All code, reference solutions, PINN checkpoints, and figure-generation scripts are released as a reproducible benchmark (GitHub: \url{https://github.com/mrmlili/aria-pinn}; Zenodo DOI: \href{https://doi.org/10.5281/zenodo.23389593}{10.5281/zenodo.23389593}).

\textbf{Fit to JAMC.} The manuscript is a computational and applied mathematics contribution that combines numerical analysis of fractional discretizations with a machine-learning solver for time-fractional PDEs. It falls squarely within the scope of JAMC.

\textbf{Honest limitations.} The manuscript openly acknowledges three limitations in Section 7.2: (i) no manufactured-solution verification of the reference solvers; (ii) no formal consistency analysis of the ABC-L1 quadrature; and (iii) small seed ensembles ($n=3$). These bound the scope of the numerical claims but do not undermine the central contributions.

\textbf{Suggested Reviewers (external, no conflict of interest).}

\textbf{1. Prof. Kai Diethelm.} Faculty of Applied Natural Sciences, Technical University of Applied Sciences W\"urzburg-Schweinfurt, Germany. Email: kai.diethelm@thws.de. Expertise: fractional calculus, numerical methods for fractional operators.

\textbf{2. Prof. Martin Stynes.} Beijing Computational Science Research Center, China. Email: m.stynes@csrc.ac.cn. Expertise: numerical analysis of fractional PDEs, graded meshes, L1 schemes.

\textbf{3. Prof. George Em Karniadakis.} Division of Applied Mathematics, Brown University, USA. Email: george\_karniadakis@brown.edu. Expertise: physics-informed neural networks, machine learning for PDEs.

\textbf{4. Prof. Dumitru Baleanu.} Department of Mathematics and Computer Science, Lebanese American University, Beirut, Lebanon. Email: dumitru.baleanu@lau.edu.lb. Expertise: fractional operators and their applications.

\textbf{Subject Classification (MSC 2020).} 65M06, 65M12, 68T07, 26A33. The manuscript sits at the intersection of fractional calculus, numerical analysis, and scientific machine learning.

The manuscript is original, has not been published elsewhere, and is not under consideration by any other journal.

Thank you for your time and consideration.

Sincerely, \\
Maryam Jalili

\end{document}
